%#############################PREAMBLE#############################################
\documentclass[a4paper,twoside,11pt,]{article}

\usepackage[spanish]{babel}
\usepackage{graphicx}
\usepackage{float}
\usepackage[skins]{tcolorbox}
\usepackage{titlepic}

\usepackage{fancyhdr}
\pagestyle{fancy}
\lhead{\emph{Machine Learning} en Geociencias}
\rhead{\thepage}
\cfoot{Programa}
\renewcommand{\headrulewidth}{0.4pt}
\renewcommand{\footrulewidth}{0.4pt}

\title {PROGRAMA  CURSO\\ \emph{Machine Learning} en Geociencias}
\author{Prof.: Edier Aristizábal\\[5ex]
\includegraphics[width=10.0cm]{G:/My Drive/FIGURAS/unal2.png}
}
\date{}

%################################BODY############################################
\begin{document}
\maketitle

\emph {versión}: \today

\section* {Introducción}
El curso \emph{Machine Learning} está orientado para estudiantes de posgrados que deseen adquirir conocimientos sobre modelos 
multivariados, para establecer patrones espaciales y análisis de inferencia estadística y predicción de una variable dependiente 
a partir de variables independentes o predictoras, con herramientas de aprendizaje automático (Machine Learning), minería de datos 
(\emph{Data Mining}) y \emph{Big Data}, bajo un ambiente de trabajo en Python.\\
El curso es teórico - práctico. Se dictarán clases teóricas con las técnicas y modelos a utilizar, pero exige que el estudiante 
seleccione una base de datos espacial para poner en practica el conocimiento adquirido en el curso. El contenido del curso incluye 
una introducción al ambiente de trabajo en Python, pero no es un curso avanzado de Python o SIG; por lo tanto, no requiere 
conocimientos profundos en dichas herramientas. Sin embargo, las personas expertas en cualquiera de estas herramientas son 
bienvenidas a aportar desde su conocimiento y en el marco del contenido del curso. El curso es una construcción conjunta de 
conocimiento por parte tanto de los estudiantes como del profesor.

\section{PROGRAMA}
El contenido del curso comprende los siguientes temas a desarrollar:

\subsection*{Introducción al curso}

\subsection {Ambiente de trabajo}

\subsection {Formulación de hipótesis}

\begin{tcolorbox}[enhanced,width=5in,center upper,  fontupper=\large\bfseries,drop shadow southwest,sharp corners]
Presentación 1 (10\%)
\end{tcolorbox}

\subsection {Toma y procesamiento de datos}

\subsection {Análisis exploratorio de datos}

\subsection {Selección de variables}

\subsection {Selección de hiperparámetros}

\subsection {Selección del modelo}

\subsection {Modelos no supervisados}

\begin{tcolorbox}[enhanced,width=5in,center upper,  fontupper=\large\bfseries,drop shadow southwest,sharp corners]
Presentación 2 -- Implementación de modelos no supervisados (30\%)
\end{tcolorbox}

\subsection {Modelos supervisados paramétricos}

\subsection {Modelos supervisados no paramétricos}

\subsection {Modelos ensamblados}

\begin{tcolorbox}[enhanced,width=5in,center upper,  fontupper=\large\bfseries,drop shadow southwest,sharp corners]
Presentación 3 -- Implementación de modelos supervisados (30\%)
\end{tcolorbox}

\subsection {\emph{Deep learning}}

\begin{tcolorbox}[enhanced,width=5in,center upper,  fontupper=\large\bfseries,drop shadow southwest,sharp corners]
Trabajo final (30\%)
\end{tcolorbox}

\section{EVALUACIÓN DEL CURSO}
El curso se evaluará con un trabajo individual donde el estudiante debe presentar los avances parciales y un informe final en formato 
de articulo. La presentación 1 corresponde a la definición del problema a resolver y la base de datos a utilizar con un valor 
del 10\% de la nota final. La presentación 2 corresponde a los resultados de la implementación de modelos no supervisados con un 
valor del 30\%. La presentación 3 corresponde a los resultados de la implementación de modelos supervisados con un valor del 30\%, 
y finalmente un trabajo escrito en formato artículo, con los resultados finales, donde se seleccione el mejor modelo y optimicen 
los resultados en los modelos de validación con un valor del 30\%. 

\end{document}
